% ================================================================================================ capítulo 6
\chapter{Modo Satélite}\label{cap:satelite}

\begin{resumen}
\textbf{En este capítulo:} cómo pedir una noche al satélite GOES (fecha, horas, intervalo y lugar), qué canales elegir,
cómo leer la serie del seeing, la tabla, los valores de cada canal y el recorte con el mapa, y qué queda guardado en la
carpeta de cada evaluación.
\end{resumen}

Antes de salir a medir, o después, para entender una medición rara, sirve saber cómo estuvo el cielo. El satélite
GOES-Este toma una imagen de todo el continente cada 10 minutos, de día y de noche, en 16 canales. El modo Satélite busca
esas imágenes en el repositorio abierto de la NOAA, lee los valores en tu sitio y, con tres canales infrarrojos, estima el
seeing de las capas medias de la atmósfera con el método de Cavazzani y colaboradores (capítulo~\ref{cap:satelite_teoria}).
Hace falta conexión a internet.

% ------------------------------------------------------------------------------------------------ formulario
\section{Pedir una noche}

\begin{capturam}{satelite}{El formulario del modo Satélite para la noche del 19 de octubre de 2026, de 18:00 a 23:00, cada
hora, con el seeing activado.}
\item \textbf{Fecha} de la noche.
\item \textbf{Desde} y \textbf{Hasta}, en hora de Colombia (UTC$-5$) o en UTC.
\item \textbf{Intervalo}: cada 30 minutos, cada 60 minutos o una sola imagen.
\item \textbf{Latitud y longitud} del sitio (norte y este positivos). \boton{OAN, Bogotá} pone las del observatorio.
\item \textbf{Recorte}: el lado del cuadro que se muestra alrededor del punto, de 10 a 1500~km.
\item \campo{Calcular el seeing}: marca los canales 9, 13 y 16.
\item \textbf{Parámetros del seeing}: longitud de onda, espesor de la capa y el cuadro en que se promedian los canales.
\item \textbf{Criterio del Sol}: qué horas entran en el cálculo.
\item \textbf{Cómo se leen los archivos}: solo la zona o los archivos completos.
\item \botonp{Buscar, leer y calcular}.
\end{capturam}

\subsection{Fecha, horas y lugar}

Con \campo{Desde} y \campo{Hasta} y un intervalo de 30 o 60 minutos se pide una serie: la línea de abajo dice cuántas
imágenes por canal son («6 imágenes por canal: de 18:00 a 23:00»). Si \campo{Hasta} es menor que \campo{Desde}, la serie
pasa la medianoche: de 22:00 a 02:00 son cinco horas. Con \campo{Una sola imagen} se lee solo la hora de \campo{Desde}.

\textbf{Satélite.} GOES-16 fue el GOES-Este hasta el 4 de abril de 2025; desde entonces lo es GOES-19. Con
\campo{Automático según la fecha}, TurbuLENS elige el que corresponde (al lado se lee cuál) y, si el archivo no está en
un repositorio, lo busca en el otro. Las imágenes de GOES-16 existen desde diciembre de 2017.

\textbf{Recorte.} Es solo para ver: el lado del cuadro que se dibuja alrededor del sitio. No puede ser menor de 10~km
(cinco píxeles de 2~km): con menos no hay imagen útil.

\subsection{Canales}

Al pasar el cursor por un canal se ve para qué sirve; la página \menu{Cómo funciona} tiene la tabla completa. Los botones
de abajo eligen grupos: \boton{Solo el 13 (nubes)}, \boton{Infrarrojos 7–16}, \boton{Todos} y \boton{Ninguno}. A la
derecha se lee cuánto se va a descargar («3 canales × 6 horas · ≈36~MB leyendo solo la zona»).

Para ver si había nubes de noche basta el canal 13, el infrarrojo limpio: las nubes altas y frías salen blancas. Con
\campo{Calcular el seeing} se marcan los canales 9, 13 y 16, que el método necesita, y no se pueden desmarcar mientras
la opción esté activa; puedes agregar otros.

\subsection{Seeing desde el satélite}

\begin{itemize}
\item \textbf{Longitud de onda} $\lambda$ (500~nm por defecto, como en el artículo del método).
\item \textbf{Espesor de la capa} $z$, en km (4~km: la capa de 4 a 8~km que mide el método).
\item \textbf{Promedio en}: el cuadro alrededor del sitio en que se promedian las temperaturas. 1°~×~1° (unos 111~km) es
el del artículo; también puedes usar 3~×~3 píxeles o solo el píxel del sitio.
\item \textbf{Criterio del Sol}: \emph{Solo de noche} (Sol a más de 18° bajo el horizonte, el recomendado),
\emph{Noche y crepúsculo} (Sol bajo el horizonte) o \emph{Todas las horas}. Las horas que no cumplen el criterio se leen
y se muestran, pero no entran en la tendencia ni en el seeing.
\end{itemize}
Al lado del criterio se ve el ángulo cenital del satélite desde tu sitio, $\theta$, y el factor
$\Lambda(\theta) = 10^{-12}/\cos\theta$. Desde Bogotá, $\theta = 5{,}62°$.

\subsection{Cómo se leen los archivos}

Cada imagen de disco completo pesa unos 25~MB por canal. Una noche de seis horas con tres canales serían 450~MB. Por eso
la opción predeterminada es \textbf{Leer solo la zona}: TurbuLENS pide al servidor únicamente los pedazos del archivo que
cubren tu recorte, unos 2~MB por canal, sin guardar nada en el disco (sección~\ref{sec:lectura_parcial}). Con
\textbf{Descargar los archivos completos} se guardan en \ruta{descargas\_goes/} dentro de la carpeta de resultados, para
volver a usarlos sin internet. Si el servidor no acepta lecturas parciales, TurbuLENS descarga el archivo completo solo.

\boton{Leer archivos ya descargados…} abre archivos \ruta{.nc} que ya tengas (también se pueden soltar en la ventana).
TurbuLENS los agrupa por la hora del barrido y arma la serie igual.

% ------------------------------------------------------------------------------------------------ resultados
\section{Los resultados}

Mientras busca y lee, la barra y el texto de abajo muestran el avance. Al terminar aparecen tres tarjetas.

\begin{capturam}{satelite_serie}{La serie de la noche: el seeing hora por hora, la función $F$ con su tendencia y la tabla.}
\item \textbf{Seeing estimado por el satélite}, hora por hora, con su mediana.
\item \textbf{$F$ y su tendencia}: la función de estabilidad de cada hora y la recta ajustada con las horas usadas. Las
horas excluidas por el Sol salen con otro símbolo.
\item \textbf{La tabla}: hora local, altura del Sol, estado (noche, crepúsculo o día), $F$, la tendencia, $|\Delta F|$,
el seeing, $r_0$ y una nota (por ejemplo «excluido por la altura del Sol»).
\item \textbf{El resumen}: cómo se calculó la tendencia y la mediana y el rango del seeing.
\end{capturam}

\textbf{Valores en tus coordenadas.} Una tabla con cada canal de la hora elegida: el valor en el píxel del sitio, el
promedio en 3~×~3 píxeles y en el cuadro, la unidad, una lectura sencilla («nube alta y fría», «probablemente
despejado»), la hora del barrido y el satélite. Al elegir una hora en la serie, esta tabla y el recorte cambian a esa hora;
al elegir un canal en esta tabla, el recorte muestra ese canal.

\begin{capturam}[0.9\linewidth]{satelite_mapa}{El recorte del canal 13 con el mapa.}
\item \textbf{El recorte}, con el sitio marcado con un círculo rojo. Al pasar el cursor se leen la latitud, la longitud
y el valor.
\item \campo{Mapa}: costas (amarillo), países (blanco) y departamentos o estados (azul), de Natural Earth.
\item \boton{Ver completo} lo abre en grande; \boton{Abrir en el visor de píxeles}, con todas las herramientas del
visor.
\item Qué se ve: canal, tamaño en píxeles y en km, y cómo leer los colores.
\end{capturam}

\begin{nota}[Con pocas horas de noche.]
La tendencia es una recta ajustada con las horas que cumplen el criterio del Sol, y para una recta hacen falta al menos
tres. Con una o dos, TurbuLENS usa su promedio y lo dice en el resumen; sin ninguna, no calcula el seeing. Si pediste una
sola imagen o un intervalo corto, amplía las horas.
\end{nota}

\begin{cuidado}[Qué mide y qué no.]
El seeing del satélite es el de la capa de 4 a 8~km. No incluye la turbulencia cerca del suelo, que en muchos sitios es
la mayor parte: por eso suele salir menor que el del DIMM. Las nubes lo invalidan. Es un método empírico, comparado con
el seeing medido en Paranal y en el Roque de los Muchachos: sirve para comparar horas y noches en un mismo sitio y para acompañar al DIMM, no para
reemplazarlo.
\end{cuidado}

\section{La carpeta de cada evaluación}\label{sec:carpeta_eval}

Para que nada se sobrescriba, cada evaluación se guarda en una subcarpeta nueva dentro de la carpeta de resultados, con
un nombre que dice qué se evaluó y cuándo:
\begin{center}
\ruta{evaluadas\_dia\_}\textbf{\ruta{20261019}}\ruta{\_}\textbf{\ruta{18-23}}\ruta{\_}\textbf{\ruta{20261126\_100953}}\\[0.3em]
{\small\color{gris}fecha de la noche \quad·\quad horas evaluadas \quad·\quad fecha y hora local en que se evaluó}
\end{center}
Con una sola imagen, las horas son solo la de inicio (\ruta{\_18\_}). Adentro van:

\begin{table}[H]
\centering\small
\begin{tabularx}{\linewidth}{@{}l Y@{}}
\toprule
Archivo & Contenido\\
\midrule
\ruta{serie\_seeing.csv} & Una fila por hora: hora local y UTC, altura del Sol, estado, si se usó, $T_9$, $T_{13}$, $T_{16}$,
$F$, la tendencia, $|\Delta F|$, $\int C_n^2\,dz$, el seeing, $r_0$ y la nota. Al final, los parámetros: sitio, $\theta$,
$\Lambda(\theta)$, $z$, $\lambda$, cuadro, criterio del Sol y tipo de tendencia.\\
\ruta{valores\_por\_canal.csv} & Cada canal y cada hora: valor en el punto, en 3~×~3 y en el cuadro.\\
\ruta{grafica\_seeing.png}, \ruta{grafica\_F\_tendencia.png} & Las dos gráficas de la serie.\\
\ruta{recortes/} & El recorte de cada canal y cada hora en FITS (con el sitio y el tamaño del píxel en la cabecera) y en PNG.\\
\ruta{LEEME.txt} & Qué se evaluó, con qué parámetros y las fórmulas.\\
\bottomrule
\end{tabularx}
\caption{Lo que queda en \ruta{evaluadas\_dia\_…/}.}
\end{table}

\section{Cómo funciona}

La segunda página del modo, \menu{Cómo funciona}, explica en tres tarjetas qué es el satélite GOES, qué mide cada uno de
sus 16 canales (con la tabla de longitudes de onda, resolución y usos) y cómo se calcula el seeing desde el espacio, con
las fórmulas. El capítulo~\ref{cap:satelite_teoria} de este manual cuenta lo mismo con más detalle.

\captura{satelite_info.png}{\menu{Cómo funciona}: el satélite, sus canales y el método.}
